% C1_identification.tex -- fix round 14 (fix14_interface), 20 Nov 2026.  Closes review item C-1 of fix11_review/REVIEW_v4b_v36_proofs.md:
% smoothness of gamma(t_1) and the identification of step 4's validated solution with the smooth CDF flow.  New file only; nothing closed
% is edited.  Not compiled (no LaTeX in WSL).  Checks: C1_sym.py -> C1_sym.json (sympy, all_pass), C1_ke.py -> C1_ke.json (arb),
% C1_num.py -> C1_num.json (float sanity check only).
\documentclass[11pt]{article}
\usepackage[margin=1in]{geometry}
\usepackage{amsmath,amssymb,amsthm,enumitem}
\newtheorem{lemma}{Lemma}[section]\newtheorem{proposition}[lemma]{Proposition}\newtheorem{theorem}[lemma]{Theorem}\newtheorem{remark}[lemma]{Remark}
\newcommand{\fn}[1]{\texttt{\detokenize{#1}}}
\newcommand{\tH}{\textup{[hand]}}\newcommand{\tS}{\textup{[sym]}}\newcommand{\tC}{\textup{[cert]}}\newcommand{\tI}{\textup{[cited]}}\newcommand{\tN}{\textup{[num]}}
\newcommand{\ip}[2]{\langle#1,#2\rangle}
\begin{document}
\title{C-1: the step-4 solution is the smooth curve diffusion flow from $\Gamma_0$,\\ and $\gamma(t_1)$ is a smooth symmetric immersion}
\date{20 November 2026 (fix round 14). \emph{Not compiled.}}
\maketitle

\noindent\textbf{Labels.} \tS: exact sympy identity in \fn{C1_sym.py} (\fn{C1_sym.json}, \fn{all_pass: true}; items Y1--Y10). \tC: ball arithmetic (\fn{python-flint arb}), script and JSON named. \tI: a result of another document, used as stated there (with its own label); for the Lean theorems this means \emph{statement read from source, not compiled}. \tH: hand argument written here, not machine-checked. \tN: float sanity check, not a proof step.

\noindent\textbf{What is closed.} Step~4 (\fn{step4_proof.tex} l.130, l.336, l.370, l.672) asserts the smoothness of $\gamma$ for $t>0$ ``by parabolic regularity'' and the continuation of the P2 solution by ``the literature fact (P1)'', both labelled [A]; and no text identifies its validated object (the P1 Galerkin solution in $C^0H^5\cap C^1L^2$, in its own gauge, then the P2 normal graph over the base) with a smooth flow. Theorem~\ref{C1:thm} below does this, with existence and uniqueness from the Lean theorems (L1), (L2), the criterion from DKS (D1), and step~4's certified bounds. It also gives (S4.1) and (S4.2) of \fn{fix10_proofs/U1_regularity.tex}, and it shows that the physical solution of step~3 is the continuation of the same flow.

%======================================================================
\section{Inputs}\label{C1:sec:in}
\paragraph{(I1) $\Gamma_0$ \tI.} Theorem~4 and Proposition ``$\Gamma_0$'' of \fn{step4_proof.tex} (\S``The initial curve''; certificate \fn{t2/gamma0_final.json}): $\theta_0(\alpha)=-\frac\pi2+\alpha+\mathrm{th}_0(\alpha)$, $\mathrm{th}_0=\sum_{n\le3762}b_n\sin n\alpha$ with exact doubles $b_n$ ($n\ne2$) and $b_2=b_2^*$ the certified unique zero of $C(b_2)=\int_0^\pi\cos(\theta_0+\frac\pi2)\,d\alpha$; $L_0$, $X_0(0)$ real. Put $X^0(\alpha)=X_0(0)+\frac{L_0}{2\pi}\int_0^\alpha e^{i\theta_0}$.
\begin{lemma}\label{C1:lem:G0}
$X^0$ is $2\pi$-periodic, real-analytic, $|X^0_\alpha|=L_0/2\pi$, and $X^0(-\alpha)=\overline{X^0(\alpha)}$. So $F_0:\mathbb S^1\to\mathbb C$, $F_0(e^{i\alpha}):=X^0(\alpha)$, is a smooth immersion with $\bar F_0=F_0\circ\iota$, $\iota(z)=\bar z$.
\end{lemma}
\begin{proof}\tH, \tS\ Y10. Since $\mathrm{th}_0$ is odd, $\theta_0(-\alpha)=-\pi-\theta_0(\alpha)$ (Y10a), so $e^{i\theta_0(-\alpha)}=-e^{-i\theta_0(\alpha)}$ (Y10b), and $X^0(-\alpha)=X_0(0)+\frac{L_0}{2\pi}\int_0^\alpha e^{-i\theta_0}=\overline{X^0(\alpha)}$. Periodicity: $\int_0^{2\pi}e^{i\theta_0}=\int_{-\pi}^{\pi}e^{i\theta_0}=\int_0^\pi(e^{i\theta_0}-e^{-i\theta_0})=2i\int_0^\pi\sin\theta_0=-2i\,C(b_2^*)=0$. $\theta_0-\alpha$ is a trigonometric polynomial, so $X^0$ is real-analytic; $|X^0_\alpha|=L_0/2\pi>0$.
\end{proof}

\paragraph{(I2) P1 and the Galerkin-limit lemma \tI.} \fn{step4_proof.tex} \S``P1'' and the appendix ``Proof of the Galerkin-limit lemma'' (ll.~236--247, 822--831). In the sine class $\mathrm{th}=\sum_{n\ge1}u_n\sin n\alpha$, with clock $d\sigma=(2\pi/L)^4dt$,
\begin{equation}\label{C1:eq:P1}
\mathrm{th}_\sigma=-\mathrm{th}''''-\mathrm{th}''+N(\mathrm{th}),\qquad N=-\bigl(v^2+\tfrac13v^3\bigr)'+(1+v)\,P\bigl(w^2-\ip{w^2}{}\bigr),\quad v=\mathrm{th}',\ w=\mathrm{th}'',
\end{equation}
$P$ the zero-mean antiderivative, $\ip{\cdot}{}$ the mean; slaved: $(\ln L)'=-\frac12\sum n^4u_n^2$, $t'=(L/2\pi)^4$, $X_0'=(L/2\pi)\sum n^3u_n$. The 96 slabs are cut into finitely many sub-intervals $I_k$ ($k=1,\dots,K$; ten per slab). On each, the lemma constructs a solution $u\in\mathcal C(I_k):=C^0(I_k;H^5)\cap C^1(I_k;L^2)$ of \eqref{C1:eq:P1} in the sine class, starting from the end value of $I_{k-1}$ (from $\mathrm{th}_0$ for $k=1$). The concatenation $u$ on $[0,\sigma_{96}]$ and the slaved $L,t,X_0$ define the \emph{P1 curve} $X^{\rm P1}(\sigma,\alpha)=X_0(\sigma)+\frac{L(\sigma)}{2\pi}\int_0^\alpha e^{i(-\pi/2+\alpha'+u)}d\alpha'$; $t(\sigma_{96})=t_{96}$, and $L(\sigma_{96})=48.293\ldots>0$.

\paragraph{(I3) Theorem P2 \tI.} \fn{step4_proof.tex} \S``P2'' (Theorem P2 and its proof, Steps 1--7; status there: constants [C], analytic steps [D], written out in \fn{LEMMA_F.md}, \fn{T7_REPORT.md}). Read precisely, Steps 1--6 and the continuity part of Step~7 are a priori statements about a given smooth CDF solution: Step~1 writes it as a normal graph $e$ over the base, Step~2 obtains $(\lambda,c,t)$ as $C^1$ functions of the label $\tau$ from the implicit function theorem (valid while $\bar m\|E\|<\frac12$; at the box $\bar m\|\hat E\|\le1.04\cdot10^{-4}$), Steps 3--6 bound the energies, and the continuity argument shows that the boxes $\mathbb B\cup\mathbb B'$ improve strictly. The only non-a-priori step is the last sentence of Step~7 (``the solution continues: this is the continuation fact (P1) of step~3, from the literature [A]''). We use P2 in this form:
\begin{enumerate}[label=(P2-\alph*),leftmargin=3em]
\item Let $\gamma$ be a smooth CDF flow on $[t_{96},T')$ with $\gamma(t_{96})$ the P1 curve at $\sigma_{96}$. Then there is a maximal $\tau_e\in(\tau_a,27.15]$ and $C^1$ functions $\lambda>0,c,t$ on $[\tau_a,\tau_e)$ with $t(\tau_a)=t_{96}$, $t(\tau)<T'$, $t_\tau=\lambda^4(1-\vartheta)$, $|\vartheta|\le3.37\cdot10^{-8}$, such that the representation $(\gamma(t(\tau))-ce_1)/\lambda=\tilde\Gamma+e\tilde\nu$, the constraints, $\sqrt{\Lambda_e}\le\bar Y=10^{-12}$, the far bounds $\sqrt{F_j}\le A_j\bar Y+B_j$ and the boxes $\mathbb B\cup\mathbb B'$ hold on $[\tau_a,\tau_e)$; and either $\tau_e=27.15$ with $t(27.15)<T'$ (then everything holds on $[\tau_a,27.15]$), or $t(\tau)\uparrow T'$ as $\tau\uparrow\tau_e$.
\end{enumerate}
(The alternative ``the modulation or a box fails while $t<T'$'' is what Step~7's strict improvement and the IFT margin exclude.) This is a reading of step~4's proof \tH, not a new estimate.

\paragraph{(I4) Lean \tI.} (L1) \fn{exists_flow_of_immersion}, (L2) \fn{flow_unique_of_immersion} of \fn{HSLflow-Whitney/CWFlow} (commit \fn{d72b7447}), as stated and instantiated in \fn{U1_regularity.tex} \S1.2: $E=\mathbb R^1$, $M=\mathbb S^1$, $V=\mathbb C$, $J=i$. \fn{IsFlow E J T Phi}: $\Phi$ jointly $C^\infty$ on $[0,T]\times\mathbb S^1$, immersions, $\partial_t\Phi=\fn{flowVel}(\Phi)$. \emph{The statements were read from source; the repository has not been built and the $n=1$ instance has not been elaborated in Lean} (review C-3). At $n=1$, $\fn{flowVel}(X)=-k_{ss}N$, $N=iT$, with no tangential part, and it commutes with $\bar\cdot$ and with reparametrisations (\fn{U1_regularity.tex} Lemma~``$n=1$'', \fn{U1_sym.py} S1--S3, \tS\ there). The lemmas ``gluing'', ``maximal smooth solution'', ``symmetry'' and ``blow-up criterion'' of \fn{U1_regularity.tex} \S2 are stated there for $F=\gamma(t_1)$; their proofs use only that $F$ is a smooth immersion with $\bar F=F\circ\iota$, so they hold verbatim for $F=F_0$ \tH.

\paragraph{(I5) DKS \tI.} (D1) (inside the proof of Thm.~3.1 of Dziuk--Kuwert--Sch\"atzle 2002, pp.~1234--35, as quoted in \fn{writeup/lit_P1/P1_CONFIRMATION.md} \S1): a smooth solution on $[0,T)$, $T<\infty$, with $\sup_{t<T}\int\kappa^2ds<\infty$ extends smoothly to $[0,T]$. \emph{Completion of review item C-4} \tH: the proof obtains the extension from uniform bounds on $\|\partial_s^m\kappa\|_{L^\infty}$, $m\ge0$, on $[0,T)$ (interpolation, with $L\ge4\pi^2/\int\kappa^2ds$ since $\int|\kappa|ds\ge2\pi$ for a closed curve); then $\partial_t\ln|\Phi_x|=-kV=kk_{ss}$ \tS\ (Y1b with $V=-k_{ss}$) is bounded, so $|\Phi_x(t)|\ge|\Phi_x(0)|e^{-CT}>0$ and the extension is an immersion at $t=T$. The DKS paper was not re-read for this sentence (the PDF could not be opened in the earlier rounds).

%======================================================================
\section{The smooth flow from $\Gamma_0$}
\begin{lemma}[maximal flow]\label{C1:lem:max}
There are $T_*\in(0,\infty]$ and a unique $\gamma:[0,T_*)\times\mathbb S^1\to\mathbb C$ with $\gamma|_{[0,T]}$ an \fn{IsFlow} from $F_0$ for every $T<T_*$, and no \fn{IsFlow} from $F_0$ on $[0,T_*]$ if $T_*<\infty$. It satisfies $\overline{\gamma(t)}=\gamma(t)\circ\iota$, and if $T_*<\infty$ then $\sup_{t<T_*}\int_{\gamma(t)}k^2ds=\infty$.
\end{lemma}
\begin{proof}\tH\ (I4) with $F=F_0$ (Lemma~\ref{C1:lem:G0}): the U-1 lemmas ``maximal smooth solution'', ``symmetry'' and ``blow-up criterion'', the last with (I5) including the C-4 completion.
\end{proof}

\noindent\emph{Gauge.} For $t<T_*$ let $p_0=1\in\mathbb S^1$ (a fixed point of $\iota$; $\alpha=0$ of $\Gamma_0$), $s(t,p)=\int_{p_0}^p|\gamma_p|$ (counter-clockwise), $L(t)=s(t,p_0^{+2\pi})$, $\alpha(t,p)=2\pi s/L$, $\theta(t,\alpha)$ the continuous lift of the tangent angle of $\gamma(t)$ at the point with normalised arclength $\alpha$, normalised by continuity in $t$ from $\theta(0,\cdot)=\theta_0$, and $\mathrm{th}^\gamma(t,\alpha)=\theta(t,\alpha)+\frac\pi2-\alpha$; $X^\gamma_0(t)=\gamma(t,p_0)$; $\sigma^\gamma(t)=\int_0^t(2\pi/L)^4$.

\begin{lemma}[the flow in step 4's gauge]\label{C1:lem:gauge}
For every $T<T_*$: (a) $\alpha$ is jointly $C^\infty$ with $\partial_p\alpha>0$, so $(t,\alpha)\mapsto\mathrm{th}^\gamma$ is jointly $C^\infty$ on $[0,T]\times\mathbb S^1$; (b) $\mathrm{th}^\gamma(t,\cdot)$ is odd and $2\pi$-periodic, i.e.\ in the sine class, and $\theta(t,0)=-\frac\pi2$; (c) as a function of $\sigma=\sigma^\gamma(t)$, $\mathrm{th}^\gamma$ solves \eqref{C1:eq:P1} pointwise, and $L$, $t$, $X_0^\gamma$ solve the three slaved equations of (I2); (d) $\mathrm{th}^\gamma\in\mathcal C([0,\sigma^\gamma(T)])$.
\end{lemma}
\begin{proof}
(a) \tH\ $|\gamma_p|$ is smooth and positive, so $s,L,\alpha$ are smooth, $L>0$, $\partial_p\alpha=2\pi|\gamma_p|/L>0$; the inverse $p(t,\alpha)$ is smooth by the inverse function theorem, and the tangent angle is a smooth function of $\gamma_p$ composed with $p(t,\alpha)$.

(b) \tH, \tS\ Y10. Differentiate $\overline{\gamma(t)}=\gamma(t)\circ\iota$ (Lemma~\ref{C1:lem:max}): $|\gamma_p|$ is $\iota$-invariant, so $s(t,\iota p)=-s(t,p)$ (mod $L$) and $\alpha\circ\iota=-\alpha$; the tangents satisfy $\overline{T(\alpha)}=-T(-\alpha)$, i.e.\ $\theta(t,-\alpha)\equiv-\pi-\theta(t,\alpha)\pmod{2\pi}$. At $\alpha=0$ this gives $\theta(t,0)\in-\frac\pi2+\pi\mathbb Z$; it is continuous in $t$ and equals $-\frac\pi2$ at $t=0$, so $\theta(t,0)=-\frac\pi2$. Both sides of $\theta(t,-\alpha)=-\pi-\theta(t,\alpha)$ are continuous in $\alpha$ and agree at $\alpha=0$, so they agree identically: $\mathrm{th}^\gamma$ is odd (Y10a). The rotation index $(\theta(t,2\pi)-\theta(t,0))/2\pi$ is an integer depending continuously on $t$, equal to $1$ at $t=0$; so $\mathrm{th}^\gamma$ is periodic.

(c) \tS\ Y1--Y7, \tH\ for the bookkeeping. With $V=-k_{ss}$ the normal speed, Y1a gives $\theta_t|_p=\partial_sV=-k_{sss}$ and Y1b gives $\partial_t|\gamma_p|=-kV|\gamma_p|$; hence $\partial_ts|_p=-\int_0^skV\,ds'$, $L_t=-\int kV\,ds$. In the gauge, Y2 expresses $k,k_s,k_{ss},k_{sss}$ through $v,w$; Y4 gives $\partial_\sigma\alpha|_p=\int_0^\alpha(1+v)w'+\alpha\ip{w^2}{}$ (with Y4b for the integrand and Y3 over a period for $L_t$); then $\theta_\sigma|_\alpha=\theta_\sigma|_p-(\partial_\sigma\alpha|_p)\theta_\alpha$ (chain rule) and Y3 (integrated from $0$, with $w(0)=0$ by oddness, Y6e) give Y5a:
\[\mathrm{th}^\gamma_\sigma=-\mathrm{th}''''-(1+v)^2w+(1+v)\int_0^\alpha\bigl(w^2-\ip{w^2}{}\bigr).\]
Y5b rewrites $(1+v)^2w=\mathrm{th}''+(v^2+v^3/3)'$. For odd $\mathrm{th}$, $w^2-\ip{w^2}{}$ is even and $\int_0^\alpha(w^2-\ip{w^2}{})$ is odd with mean $0$, hence equals $P(w^2-\ip{w^2}{})$ (Y6b--d, checked on a symbolic odd trigonometric polynomial of degree 3; the general case is the same two-line parity argument \tH). This is \eqref{C1:eq:P1}. Slaved equations: $\frac{d}{d\sigma}\ln L=-\frac1{2\pi}\int w^2=-\frac12\sum n^4u_n^2$ (Y7a, Y6a; the Parseval step for infinitely many modes is \tH); $X_0^\gamma$ is the image of the fixed parameter $p_0$, so $\partial_tX^\gamma_0=V(p_0)N(p_0)$ with $N(p_0)=ie^{-i\pi/2}=1$ (Y7c) and $\frac{d}{d\sigma}X^\gamma_0=-(L/2\pi)\mathrm{th}'''(0)=(L/2\pi)\sum n^3u_n$ (Y7b); $t'=(L/2\pi)^4$ by definition of $\sigma$.

(d) \tH\ $\sigma^\gamma$ is a $C^1$ diffeomorphism of $[0,T]$ onto $[0,\sigma^\gamma(T)]$, and a jointly smooth function on a compact set is in $C^0H^5\cap C^1L^2$.
\end{proof}
\noindent\tN\ \fn{C1_num.py}: one explicit normal step of a closed symmetric test curve, re-gauged numerically, reproduces \eqref{C1:eq:P1} on the low modes with relative defect $3.0\cdot10^{-6}$ ($dt=10^{-6}$) and $3.1\cdot10^{-7}$ ($dt=10^{-7}$), and $(\ln L)'$ with $3.6\cdot10^{-5}$ ($dt=10^{-7}$). This is independent of the derivation above (it never uses Y4/Y5) and is only a cross-check.

\begin{lemma}[uniqueness in $\mathcal C$, without a box]\label{C1:lem:uniq}
Let $u,\tilde u\in C^0([0,S];H^5)\cap C^1([0,S];L^2)$ be odd, solve \eqref{C1:eq:P1} in $L^2$ for every $\sigma$, and $u(0)=\tilde u(0)$. Then $u\equiv\tilde u$.
\end{lemma}
\begin{proof}\tH\ (the argument of step~6 of the Galerkin-limit proof, \fn{step4_proof.tex} l.829, written for the whole class: the lemma's statement claims uniqueness only among solutions staying in its isolating set $W$, while its proof does not use $W$.) Let $\delta=u-\tilde u$, $f=\|\delta\|^2$ ($L^2(0,2\pi)$), so $f\in C^1$, $f'=2\ip{\delta_\sigma}{\delta}$. Let $M:=\sup_{[0,S]}(\|v\|_\infty+\|\tilde v\|_\infty+\|w\|+\|\tilde w\|)<\infty$ ($H^5\subset C^4$, continuity on a compact interval). Periodic integration by parts: $\ip{-(\partial^4+\partial^2)\delta}{\delta}=-\|\delta''\|^2+\|\delta'\|^2$. With $g(v)=v^2+v^3/3$, $|g(v)-g(\tilde v)|\le(2M+M^2)|\delta'|$, so $|\ip{-(g(v)-g(\tilde v))'}{\delta}|=|\ip{g(v)-g(\tilde v)}{\delta'}|\le(2M+M^2)\|\delta'\|^2$. Next $(1+v)P(w^2)-(1+\tilde v)P(\tilde w^2)=\delta'P(w^2)+(1+\tilde v)P(\delta''(w+\tilde w))$ ($P$ kills constants), $\|Pf\|_\infty\le4\|f\|_{L^1}$, so this pairs with $\delta$ to at most $4\sqrt{2\pi}M^2\|\delta'\|\|\delta\|+8\sqrt{2\pi}(1+M)M\|\delta''\|\|\delta\|$. With $\|\delta'\|^2\le\|\delta\|\|\delta''\|$ and Young, $\frac12f'\le-\frac12\|\delta''\|^2+C(M)f$; Gr\"onwall and $f(0)=0$ give $f\equiv0$.
\end{proof}

%======================================================================
\section{Identification on P1 ($0\le t\le t_{96}$)}
\begin{lemma}[curvature bound on P1]\label{C1:lem:kP1}
$\sup_{\sigma\le\sigma_{96}}\int_{X^{\rm P1}(\sigma)}k^2ds<\infty$.
\end{lemma}
\begin{proof}\tH\ $\int k^2ds=\frac{2\pi}L\int_0^{2\pi}(1+v)^2d\alpha$ (Y2a), $L$ is nonincreasing (slaved equation) so $L\ge L(\sigma_{96})>0$, and $\sup_{[0,\sigma_{96}]}\|v\|_\infty<\infty$ because $u\in C^0(I_k;H^5)$ on each of the finitely many $I_k$.
\end{proof}

\begin{proposition}[identification on P1]\label{C1:prop:P1}
$T_*>t_{96}$, $\sigma^\gamma(t_{96})=\sigma_{96}$, and for $t\in[0,t_{96}]$: $\mathrm{th}^\gamma(t)=u(\sigma^\gamma(t))$, $L$, $X_0$ agree, and $\gamma(t)$ is the P1 curve up to the reparametrisation $p\mapsto\alpha(t,p)$. In particular the P1 enclosures are enclosures of the smooth flow.
\end{proposition}
\begin{proof}\tH\ For $t'<T_*$ let $S'=\min(\sigma^\gamma(t'),\sigma_{96})$. By Lemma~\ref{C1:lem:gauge}, $\mathrm{th}^\gamma$ is in the class of Lemma~\ref{C1:lem:uniq} on $[0,S']$ and starts from $\mathrm{th}_0$; applying Lemma~\ref{C1:lem:uniq} on $I_1\cap[0,S'],I_2\cap[0,S'],\dots$ in turn gives $\mathrm{th}^\gamma=u$ on $[0,S']$. Then $L,t,X_0$ agree on $[0,S']$ (same ODEs, same data), so the curves agree. Case A: $\sigma^\gamma(t')>\sigma_{96}$ for some $t'<T_*$; then the flow time at $\sigma_{96}$ is $t(\sigma_{96})=t_{96}<t'<T_*$, which is the claim. Case B: $\sigma^\gamma(t)\le\sigma_{96}$ for all $t<T_*$; then every $\gamma(t)$, $t<T_*$, is a P1 curve, so $\int k^2ds$ is bounded (Lemma~\ref{C1:lem:kP1}) and $T_*=\infty$ (Lemma~\ref{C1:lem:max}); but $L\le L_0$ gives $\sigma^\gamma(t)\ge(2\pi/L_0)^4t\to\infty$, a contradiction.
\end{proof}

%======================================================================
\section{Identification on P2 ($t_{96}\le t\le t_1$)}
\paragraph{What P2 certifies.} Theorem~P2 is a statement about a smooth (classical) CDF solution, used only through geometric quantities: its normal velocity (Step~1 uses $\ip{\Gamma_\tau}{\nu}$ only), its normal graph $e$ over the base, and energies of $e$ (Step~1: $e\in W^{7,\infty}$ in $x$, identities a.e.\ in $\tau$). These are independent of the parametrisation, so the Lean flow, whose velocity is purely normal, is directly in P2's class; no gauge change is needed (unlike P1). Its initial datum is the P1 state at $t_{96}$ (switch lemma), which is $\gamma(t_{96})$ by Proposition~\ref{C1:prop:P1}.

\begin{lemma}[the normal-graph factor]\label{C1:lem:ke}
On the boxes $\mathbb B\cup\mathbb B'$ of P2, for every slab $2\ldots180$ and every stored cell of the corrected base, $|\tilde\kappa\,e|\le1.743\cdot10^{-3}$; hence $a:=1-\tilde\kappa e\ge\frac12$ and $J:=(a^2+e'^2)^{1/2}\ge\frac12$.
\end{lemma}
\begin{proof}\tC\ \fn{C1_ke.py} $\to$ \fn{C1_ke.json} (\fn{sup_ke} $=1.7428\cdot10^{-3}$ at slab 180, $x\approx3.6\cdot10^5$; box $|e|\le32.3$ there): hulls of $\tilde\kappa$ from \fn{p2/t6c/jets} (widened outward by 2 ulps), $|e|\le10^{-5}$ for $y<28$, $|e|\le1.7E(x)$ for $y\ge28$ with $\sqrt{F_0}\le a_F\bar Y+b_F$ per slab (\fn{amp_A1_w180.json}) at $\bar Y=1.7\cdot10^{-12}$ (conservative). All 180 input files are hashed in the JSON.
\end{proof}

\begin{lemma}[curvature bound on P2's interval]\label{C1:lem:kP2}
In the situation of (P2-a), $\sup_{\tau\in[\tau_a,\tau_e)}\int_{\gamma(t(\tau))}k^2ds<\infty$.
\end{lemma}
\begin{proof}\tH, \tS\ Y8. The curve is $ce_1+\lambda(\tilde\Gamma+e\tilde\nu)$ on the half $x\in[0,\tilde\ell]$ and its mirror image. With $x$ the base arclength, $ds_{\rm phys}=\lambda J\,dx$, $k_{\rm phys}=\kappa/\lambda$, and $\kappa=(ae''-e'a'+\tilde\kappa(a^2+e'^2))/J^3$ (Y8b), $a'=-\tilde\kappa'e-\tilde\kappa e'$. Hence $\int k^2_{\rm phys}ds_{\rm phys}=\frac2\lambda\int_0^{\tilde\ell}\kappa^2J\,dx\le\frac{C}{\lambda}\int_0^{\tilde\ell}(1+e''^2)dx$ with $C$ depending on $\sup(|e|,|e'|)$ (boxes), $\sup(|\tilde\kappa|,|\tilde\kappa'|)$ (the base is $C^4$ with enclosed jets on the compact label range) and $J\ge\frac12$ (Lemma~\ref{C1:lem:ke}). On $y\le29$ ($\chi=1$), $\int e''^2\le\rho(31.25)^{-1}\|(\chi e)''\|_\rho^2\le\rho(31.25)^{-1}\bar Y^2/\varepsilon$ ($\Lambda_e\ge\varepsilon B^2$). On $y\ge28$ the box gives $x|e''|\le b_1$, and on $x\ge10.25$, $y<28$ it gives $x|e''|\le\eta_F$, so $\int e''^2<\infty$ uniformly. $\tilde\ell\le3.7\cdot10^5$. Finally $\lambda_\tau/\lambda=a_{\rm mod}-\frac14$ with $|a_{\rm mod}-\tilde a|\le|x_1|+|\vartheta c|$ bounded (Step~2: $|x|\le C_xA+x_r+x_{\rm NL}$) and $\tilde a$ bounded, so $\lambda\ge\lambda(\tau_a)e^{-C(27.15-\tau_a)}>0$.
\end{proof}

\begin{proposition}[identification on P2]\label{C1:prop:P2}
$T_*>t(27.15)\ge t_1$, and on $[t_{96},t(27.15)]$ the conclusions of Theorem~P2 hold for the smooth flow $\gamma$ of Lemma~\ref{C1:lem:max}.
\end{proposition}
\begin{proof}\tH\ Apply (P2-a) with $T'=T_*$. If $t(\tau)\uparrow T_*$ as $\tau\uparrow\tau_e$, then $T_*\le t_{96}+\int_{\tau_a}^{27.15}\lambda^4<\infty$, and $\int k^2ds$ is bounded on $[0,t_{96}]$ (Lemma~\ref{C1:lem:kP1} and Proposition~\ref{C1:prop:P1}) and on $[t_{96},T_*)$ (Lemma~\ref{C1:lem:kP2}), contradicting Lemma~\ref{C1:lem:max}. So $\tau_e=27.15$ and $t(27.15)<T_*$. $t$ is increasing ($t_\tau=\lambda^4(1-\vartheta)>0$), and $t_1=t(\tau^*)$ with $\tau^*\le27.15$ (slab 180).
\end{proof}

%======================================================================
\section{The theorem}
\begin{theorem}[C-1]\label{C1:thm}
Let $\gamma$ be the maximal smooth CDF flow from $\Gamma_0$ given by the Lean theorems (Lemma~\ref{C1:lem:max}). Then $T_*>t(27.15)\ge t_1$ for every good time $\tau^*$ of slab 180, and:
\begin{enumerate}[label=(\roman*),leftmargin=2em]
\item $\gamma$ is jointly $C^\infty$ on $[0,t(27.15)]\times\mathbb S^1$, each $\gamma(t)$ is a closed immersed curve, symmetric: $\overline{\gamma(t)}=\gamma(t)\circ\iota$; this includes $t=0$ and $t=t_1$;
\item on $[0,t_{96}]$, step~4's P1 solution is the tangent angle of $\gamma$ in normalised-arclength gauge, with the same $L$, $X_0$ and clock (Proposition~\ref{C1:prop:P1});
\item on $[t_{96},t(27.15)]$, step~4's P2 decomposition is a decomposition of $\gamma$ (Proposition~\ref{C1:prop:P2});
\item hence every certified hand-off item of Theorem~4(c), (H1)--(H5), is a statement about $\gamma(t_1)$, and (H6) holds: $\gamma(t_1)$ is a smooth closed immersed curve. In the notation of \fn{U1_regularity.tex}, (S4.1) and (S4.2) hold, the latter with $\iota(z)=\bar z$;
\item the maximal flow of \fn{U1_regularity.tex} Lemma ``maximal smooth solution'' (from $F=\gamma(t_1)$) is $t\mapsto\gamma(t_1+t)$, $t\in[0,T_*-t_1)$, i.e.\ its $T_{\max}=T_*-t_1$: step~3 continues the same flow.
\end{enumerate}
\end{theorem}
\begin{proof}\tH\ (i)--(iii) are Lemma~\ref{C1:lem:max} and Propositions~\ref{C1:prop:P1}, \ref{C1:prop:P2} ($\gamma|_{[0,T]}$ is an \fn{IsFlow} for $T<T_*$, in particular on $[0,t(27.15)]$). (iv): step~4's hand-off quantities are functions of the true curve at $t_1$, which by (iii) is $\gamma(t_1)$. (v): $t\mapsto\gamma(t_1+t)$ restricted to $[0,T]$, $T<T_*-t_1$, is an \fn{IsFlow} from $\gamma(t_1)$ (the clauses are invariant under a time shift); by (L2) and the gluing lemma both maximal flows coincide where defined and neither can be longer than the other.
\end{proof}
\begin{remark}
This replaces step~4's sentences ``smooth for $t>0$ by parabolic regularity'' (l.672, [A]) and ``continuation uses the literature fact (P1)'' (ll.130, 336, 370, [A]) by Lemma~\ref{C1:lem:max} and Propositions~\ref{C1:prop:P1}, \ref{C1:prop:P2}. It does not use step~4's (H6) or step~3's (P1$'$) as an input, so there is no circularity: (P1$'$) for step~3 is (v) together with \fn{U1_regularity.tex}. Smoothness holds on the closed interval from $t=0$, not only for $t>0$, because $\Gamma_0$ is a trigonometric curve; no parabolic smoothing is used.
\end{remark}

%======================================================================
\section{Gaps, stated plainly}
\begin{enumerate}[leftmargin=*]
\item \textbf{Lean.} (L1), (L2) are read from source at \fn{d72b7447}; not compiled; the $n=1$ instance is not elaborated (review C-3, owed by Micah: \fn{lake build} and a short instantiation file). The closure of \fn{IsFlow} under $\bar\cdot\circ\iota$, the gluing and the time shift are mathematical arguments, not Lean lemmas.
\item \textbf{DKS (D1)} is a step inside a proof, read through \fn{P1_CONFIRMATION.md}; the C-4 completion (uniform $\|\partial_s^m\kappa\|_\infty$ in that proof) is my reading of the structure of DKS's argument and has not been checked against the paper.
\item \textbf{P2 as an a priori theorem} (I3) is a reading of step~4's proof; P2's analytic steps keep step~4's status ([D], written out in \fn{LEMMA_F.md}, \fn{T7_REPORT.md}). Lemma~\ref{C1:lem:kP2} uses P2's boxes and Step~2's bound on $|x|$ qualitatively; only the factor $a\ge\frac12$ is a new certificate.
\item \textbf{Galerkin existence} (steps 1--5 of step~4's appendix) is used as stated there (reviewed in step~4). Lemma~\ref{C1:lem:uniq} re-proves only uniqueness.
\item \textbf{Parseval and parity for infinitely many modes} (Lemma~\ref{C1:lem:gauge}(c)) are \tH; the sympy checks Y6, Y7 are on symbolic trigonometric polynomials of degree 3.
\end{enumerate}
\end{document}
